\documentclass[11pt,a4paper]{article}
\usepackage[margin=1in]{geometry}
\usepackage{hyperref}
\usepackage{enumitem}
\usepackage{graphicx}
\usepackage{xcolor}

% -----------------------------------------------------
% bvraghav-tiet-ucs503p-article.sty
% -----------------------------------------------------

% Variable: Lab Instructor
\makeatletter \newcommand{\BvrSetLabInstructor}[1]{%
  \gdef\Bvr@LabInstructor{#1}}
% default
\newcommand{\Bvr@LabInstructor}{Dr. Raghav %
  B. Venkataramaiyer}
% public accessor
\newcommand{\BvrLabInstructorValue}{\Bvr@LabInstructor}
\makeatother

% Add subtitle
\usepackage{titling}
% -- Set title bold
\pretitle{\begin{center}\bfseries\fontsize{18bp}{18bp}\selectfont}
\makeatletter
\newcommand{\BvrSubtitle}[1]{%
  \posttitle{%
    \par\end{center}
    \begin{center}\large#1\end{center}
    \vskip 0.5em}%
}
\makeatother

% Hints in the section title(s)
\newcommand{\BvrHint}[1]{%
  {\normalfont\normalsize\itshape\hfill #1}}

% -----------------------------------------------------

\title{FOCUS.OS}

\BvrSetLabInstructor{Paramveer Sidhu}

\BvrSubtitle{%
  Project Proposal for UCS503  \\
  Submitted to: \BvrLabInstructorValue \\ \bfseries
  Thapar Institute of Engineering and Technology%
}

\author{
  Deeksha Tanwar%
  \thanks{Roll No: \texttt{1024170384}, %
    Email: \texttt{dtanwar\_be24@thapar.edu}}
  \and
  Daksh Sharma%
  \thanks{Roll No: \texttt{1024170385}, %
    Email: \texttt{dsharma1\_be24@thapar.edu}}
}

\date{\today}

\begin{document}
\maketitle

\section{Higher-order goal%
  \BvrHint{\ldots secondary framing}}

This project aims to improve everyday productivity and task management for
people who experience difficulties with planning, task initiation,
prioritization, maintaining routines, and managing multiple responsibilities.
The system is particularly intended to support users with ADHD and similar
executive-function challenges, while remaining useful to a broader user
population.

The higher-order goal is to reduce the cognitive effort required to decide
what to do next and to help users make consistent progress on their
responsibilities.

\section{Time-to-value %
\BvrHint{\ldots primary heuristic}}

\begin{itemize}[leftmargin=0.7cm]
    \item We will deliver a usable task-management and planning workflow
    within the first iteration.
    
    \item We will prioritize rapid user feedback through weekly iterations,
    allowing the planning and interaction model to evolve based on observed
    user behaviour and feedback.
    
    \item The initial system will focus on task capture, prioritization, and
    a simple ``What should I do now?'' workflow before introducing more
    advanced personalization and AI-based features.
    
    \item CI/CD will be introduced early to automate testing, building, and
    deployment of incremental versions.
\end{itemize}

\section{Problem Statement}

People often manage their responsibilities across multiple disconnected
systems such as calendars, reminder applications, notes, habit trackers,
and paper lists. Even when users know what they need to accomplish, deciding
what to work on next and maintaining a realistic plan can become difficult.

This problem can be particularly significant for people experiencing
executive-function difficulties, including difficulties commonly associated
with ADHD.

Common pain points include:

\begin{itemize}[leftmargin=0.7cm]
    \item \textbf{Task overload:} users may have many responsibilities but
    struggle to determine which task should be handled first.
    
    \item \textbf{Task initiation:} knowing that something needs to be done
    does not necessarily translate into starting it.
    
    \item \textbf{Large or vague tasks:} tasks such as ``complete project''
    or ``study for exam'' can be difficult to start without smaller,
    actionable steps.
    
    \item \textbf{Rigid schedules:} missing one task can make a manually
    planned schedule unrealistic for the rest of the day.
    
    \item \textbf{Forgetting recurring responsibilities:} habits, routines,
    and repeated tasks may be difficult to maintain consistently.
    
    \item \textbf{Decision fatigue:} users may spend significant time deciding
    what to work on instead of actually working.
\end{itemize}

\textbf{Why this matters:} Existing productivity tools often provide task
lists, calendars, reminders, or habit tracking, but the user may still be
required to decide how to organize those tasks and what to do at a particular
moment.

\section{Proposed Solution %
\BvrHint{\ldots Software Proposition}}

\subsection{Overview}

Build a \textbf{web-based adaptive productivity and task-planning system}
that allows users to put their responsibilities, routines, and recurring
activities into one system. The system continuously organizes these inputs
and helps the user determine the most appropriate next action.

The main components are:

\begin{itemize}[leftmargin=0.7cm]
    \item \textbf{Task management:} users can quickly add tasks along with
    deadlines, priority, and estimated duration.
    
    \item \textbf{Smart prioritization:} the system ranks tasks using factors
    such as urgency, priority, estimated duration, and available time.
    
    \item \textbf{Dynamic planning:} the system adapts the plan when tasks are
    delayed, skipped, completed early, or when the user's available time
    changes.
    
    \item \textbf{Task breakdown:} larger tasks can be divided into smaller,
    actionable steps.
    
    \item \textbf{Routines and habits:} users can define recurring activities
    such as morning routines, exercise, reading, or other regular tasks.
    
    \item \textbf{Focus mode:} users can start a focused work session for the
    currently recommended task using a timer and progress tracking.
    
    \item \textbf{Gamification:} completion of tasks and focus sessions can
    provide points, levels, or user-defined rewards to encourage consistency.
\end{itemize}

\subsection{Core workflow}

\begin{enumerate}[leftmargin=0.7cm]
    \item User quickly captures the tasks and responsibilities they need to
    complete.
    
    \item The system stores relevant information such as deadlines,
    priorities, estimated duration, and recurrence.
    
    \item The system evaluates the user's available time and current tasks
    and determines which task should be addressed next.
    
    \item The user starts the recommended task through Focus Mode.
    
    \item After the task is completed, postponed, or interrupted, the system
    updates the remaining plan.
    
    \item The system continuously adapts the schedule instead of requiring
    the user to manually rebuild the entire day's plan.
\end{enumerate}

\subsection{Example interaction}

A user enters:

\begin{quote}
``Complete Software Engineering project by Friday.''
\end{quote}

The system may break the task into smaller steps such as:

\begin{itemize}[leftmargin=0.7cm]
    \item Review project requirements
    \item Define functional requirements
    \item Design system architecture
    \item Design database
    \item Implement core functionality
    \item Test the system
    \item Prepare deployment and documentation
\end{itemize}

If the user later has only 30 minutes available, the system may recommend a
small actionable step rather than presenting the entire project as one task.

\section{Solution Approach %
\BvrHint{\ldots Engineering focus}}

This is an engineering course project, so the focus will be on building a
reliable adaptive planning workflow rather than developing a novel AI or
clinical system.

\subsection{Smart planning}

The initial prioritization engine will use measurable factors such as:

\begin{itemize}[leftmargin=0.7cm]
    \item task deadline,
    \item user-defined priority,
    \item estimated duration,
    \item available time,
    \item task dependencies, where applicable,
    \item recurring commitments.
\end{itemize}

The system will provide an understandable reason for its recommendation,
rather than making the prioritization decision completely opaque.

For example:

\begin{quote}
\textbf{Work on DBMS assignment now}

Due tomorrow, approximately 40 minutes remaining, and
you currently have 50 minutes available.
\end{quote}

\subsection{Adaptive scheduling}

The system will maintain a dynamic plan rather than treating the schedule
as fixed. When a task is missed or interrupted, remaining tasks can be
reordered based on their current urgency, priority, and available time.

\subsection{AI-assisted functionality}

AI may be used as a supporting component for tasks such as:

\begin{itemize}[leftmargin=0.7cm]
    \item converting vague tasks into actionable subtasks,
    \item suggesting task duration,
    \item assisting with schedule creation,
    \item suggesting an appropriate next action,
    \item adapting plans based on user-provided context.
\end{itemize}

AI functionality will be introduced only where it provides measurable
value and will not be treated as the primary novelty of the project.

\subsection{Web architecture}

A typical 3-tier architecture will be used:

\begin{itemize}[leftmargin=0.7cm]
    \item \textbf{Frontend:} task dashboard, planner, routines, focus mode,
    and progress views.
    
    \item \textbf{Backend API:} authentication, task management,
    prioritization, scheduling, habit/routine logic, and progress tracking.
    
    \item \textbf{Database:} users, tasks, deadlines, routines, habits,
    focus sessions, task history, and reward information.
\end{itemize}

\subsection{CI/CD and fast delivery enabler}

CI/CD will be used to:

\begin{itemize}[leftmargin=0.7cm]
    \item automatically build and test the application,
    \item run unit and integration tests on changes,
    \item produce repeatable deployment artifacts,
    \item enable frequent deployment of incremental features.
\end{itemize}

\section{Evaluation Criterion %
\BvrHint{\ldots measurable and attributable}}

The effectiveness of the system will be evaluated using measurable
user-facing and system-level metrics.

\subsection{Primary evaluation metrics}

\begin{itemize}[leftmargin=0.7cm]
    \item \textbf{Task initiation time:} time between a recommended task being
    presented and the user starting the task.
    
    \item \textbf{Planning effort:} time required for users to decide what
    task to work on compared with a manual planning approach.
    
    \item \textbf{Task completion rate:} percentage of planned tasks completed
    within the intended time period.
\end{itemize}

\subsection{Secondary evaluation metrics}

\begin{itemize}[leftmargin=0.7cm]
    \item \textbf{Schedule recovery:} percentage of delayed schedules that
    successfully generate a revised plan without manual replanning.
    
    \item \textbf{Recommendation usefulness:} user rating of whether the
    recommended next task was appropriate.
    
    \item \textbf{Task breakdown usefulness:} user rating of whether generated
    subtasks made large tasks easier to start.
    
    \item \textbf{Focus-session completion:} percentage of started focus
    sessions completed successfully.
    
    \item \textbf{User satisfaction:} simple 1--5 rating collected during
    pilot testing.
\end{itemize}

\subsection{Pilot validation plan %
\BvrHint{\ldots high-level}}

\begin{itemize}[leftmargin=0.7cm]
    \item Recruit a small group of users for initial usability testing.
    
    \item Compare manual planning with the proposed system for representative
    task-management scenarios.
    
    \item Collect task initiation time, completion rate, planning effort,
    and user feedback.
    
    \item Use the results to refine the prioritization and interaction
    workflow in subsequent iterations.
\end{itemize}

\section{Scalability %
\BvrHint{\ldots with foresight}}

The proposition should scale in both theoretical and practical contexts.

\begin{enumerate}[leftmargin=0.7cm]
    \item \textbf{Scalable (theoretically):} The system should support an
    increasing number of users and tasks through:
    
    \begin{itemize}[leftmargin=0.7cm]
        \item stateless backend APIs,
        \item indexed database queries,
        \item pagination for large task histories,
        \item separation of frontend and backend services.
    \end{itemize}
    
    \item \textbf{Operationally deployable (practically):} The application
    should be deployable using common web hosting, managed databases, and
    automated CI/CD pipelines.
\end{enumerate}

\section{Engine availability heuristic}

A mature set of engineering tools and services is available for implementing
the proposed system:

\begin{itemize}[leftmargin=0.7cm]
    \item Web framework for the frontend and backend API.
    \item Relational or document database for persistent user data.
    \item Authentication and authorization mechanisms.
    \item Scheduling and notification services.
    \item AI/LLM APIs for optional task decomposition and planning assistance.
    \item Test frameworks for unit and integration testing.
    \item CI runner and deployment platform.
\end{itemize}

These mature components allow the project to focus on software engineering,
workflow design, system reliability, user feedback, and measurable outcomes
rather than research into new AI algorithms.

\section{Project scope and deliverables %
\BvrHint{\ldots iteration-friendly}}

\subsection{Initial deliverable %
\BvrHint{\ldots first iteration}}

\begin{itemize}[leftmargin=0.7cm]
    \item User registration and basic profile
    \item Quick task creation
    \item Task deadlines, priorities, and estimated duration
    \item Basic task prioritization
    \item ``What should I do now?'' recommendation
    \item Focus timer
    \item Basic task completion tracking
    \item CI: build + backend unit tests + linting
    \item CD to staging
\end{itemize}

\subsection{Subsequent deliverables}

\begin{itemize}[leftmargin=0.7cm]
    \item Dynamic schedule adjustment after missed or interrupted tasks
    \item Task decomposition into smaller actionable steps
    \item Recurring routines and habits
    \item AI-assisted planning and task breakdown
    \item Gamification and user-defined rewards
    \item Improved personalization using historical task behaviour
    \item Notification and quick-access mechanisms
    \item Calendar integration, if feasible within the project timeline
\end{itemize}

\section{Risks and mitigations}

\begin{itemize}[leftmargin=0.7cm]
    \item \textbf{App avoidance:} provide quick task capture, minimal interaction,
    and optional notifications/widgets so users do not need to repeatedly
    navigate through the application.
    
    \item \textbf{Information overload:} keep the primary interface focused on
    the single most useful question: ``What should I do now?''
    
    \item \textbf{Incorrect prioritization:} provide transparent reasons for
    recommendations and allow users to override them.
    
    \item \textbf{Inaccurate time estimates:} allow users to modify estimates
    and use historical completion data for future suggestions.
    
    \item \textbf{Repeated procrastination:} identify frequently postponed
    tasks and offer smaller actionable steps or rescheduling.
    
    \item \textbf{Schedule disruption:} dynamically recalculate remaining
    tasks when the user falls behind or available time changes.
    
    \item \textbf{Notification fatigue:} prioritize important notifications
    rather than generating reminders for every activity.
    
    \item \textbf{Feature overload:} develop the core task-planning workflow
    first and treat habits, gamification, and advanced AI functionality as
    subsequent iterations.
    
    \item \textbf{Privacy:} minimize collection of personal information and
    provide appropriate controls for task, routine, and behavioural data.
    
    \item \textbf{ADHD-related claims:} position the system as a productivity
    and planning tool designed with executive-function challenges in mind,
    rather than as a medical or therapeutic product.
\end{itemize}

\section{Summary}

This project proposes a web-based adaptive task-planning and focus-support
system designed around the difficulties of managing multiple responsibilities,
initiating tasks, prioritizing work, and recovering from disrupted schedules.

Rather than functioning only as a conventional task list, the system will
continuously organize user inputs and provide a context-aware recommendation
for what the user should work on next. The system will progressively
incorporate task breakdown, routines, focus sessions, adaptive scheduling,
AI assistance, and gamification.

The project meets the course criteria by emphasizing rapid incremental
delivery, CI/CD-supported development, mature engineering components,
scalability, and measurable evaluation through task initiation time,
completion rate, planning effort, and schedule recovery.

\end{document}
